\documentclass[10pt]{article}
\usepackage[margin=20mm]{geometry}
\usepackage{amsmath,amssymb,mathtools,fontspec,unicode-math,xurl,hyperref}
\IfFontExistsTF{DejaVu Serif}{\setmainfont{DejaVu Serif}}{\setmainfont{Latin Modern Roman}}
\setmathfont{latinmodern-math.otf}
\setlength{\parindent}{0pt}\setlength{\parskip}{6pt}\sloppy
\title{A missing two-sided hypothesis in a good-\(\lambda\) moment estimate}
\author{AI technical review}
\date{}
\begin{document}
\maketitle
Abstract: This report checks the testing conditions and the passage from one-sided tail estimates to the full \(L^p\) bound in Theorem 3.6 of the acquired author manuscript. A finite commutative martingale satisfies every stated hypothesis at every positive scale but violates the asserted bound at \(p=4\). The unsupported step is the extension of a one-sided test to the opposite sign. The finding concerns arXiv:1805.07057v2 only; correspondence with the later publisher text has not been established. Scientific adjudication remains pending.\par
\section*{Identity, version, and coverage}
The assigned paper is Noncommutative good-\(\lambda\) inequalities, by Yong Jiao, Adam Osękowski, and Lian Wu, frozen identifier M25-0271 and journal DOI https://doi.org/10.4171/JEMS/1644. The actual reviewed document is arXiv:1805.07057v2, dated 19 November 2018, 52 pages [source]. All page references below are one-based PDF pages; these coincide with the printed page numbers in this source.\par
The publisher record gives publication on 7 May 2025, submission on 7 September 2023, and acceptance on 5 August 2024 [publisher]. The arXiv landing page separately lists a 2024 journal-reference year. These facts do not establish textual equivalence. The publisher full text was not available in the supplied source set; no current-journal defect is asserted.\par
The review read the introduction and preliminaries, followed Definitions 3.2–3.3 and Theorems 3.4–3.6 through their displayed proofs, and checked the sign dependence of the application in Theorem 4.1 (pp. 1–20). Theorem 4.2 was read only through its opening construction on p. 20. The remaining applications in Sections 4–6 and their external proof dependencies were not independently verified. Every supplied PDF page was parsed, which is distinct from this narrower scientific coverage. Source images for pp. 1, 5, and 9–20 were actually viewed.\par
\section*{The tested assertion}
Equation reference: test1, test2, asserted. Throughout Section 3 the martingale \(y=(y_n)_{n=0}^N\) is self-adjoint; it is not assumed positive. The differences include \(dy_0=y_0\) (p. 5). For a positive threshold \(\mu\), write the source’s stopping projections as \[R_{-1}^{\mu}=I,\qquad
 R_n^\mu=R_{n-1}^\mu
 \mathbf 1_{(-\infty,\mu)}(R_{n-1}^\mu y_nR_{n-1}^\mu).\] At threshold one, Definition 3.2 (p. 10) requires \[\begin{aligned}
 \sum_{n=0}^N\sum_{k=n+1}^N
 \tau((R_{n-1}-R_n)dy_kR_{n-1}dy_k(R_{n-1}-R_n))
 &\le \tau((I-R_N)x_N^2), \\
 \tau(Pdy_k^2P)&\le \tau(Pz_N^2P)
 \quad(0\le k\le N,\ P\in\mathcal M_k\text{ a projection}).\end{aligned}\] Theorem 3.6 (p. 17, (3.25)) assumes these conditions for \((x_N/\mu,y/\mu,z_N/\mu)\) for every \(\mu>0\), with \(x_N,z_N\in L^p\), and asserts, for \(2<p<\infty\), \[
 \|y_N\|_p\le
 \frac{12p}{\sqrt{1-(1+1/p)^{2-p}}}
 \bigl(\|x_N\|_p^2+\|z_N\|_p^2\bigr)^{1/2}.\] Testing the upper crossings of \(y\) is intrinsically one-sided. The stronger conditions in Definition 3.3 are different assumptions; they involve squared differences and are sign invariant.\par
\section*{Substantive finding}
Problem 1 (The full moment bound does not follow from the stated one-sided tests).  Theorem 3.6 is false under its printed hypotheses. In its proof, the negative-side estimate (3.29) on p. 18 does not follow from the positive-side estimate (3.28).\par
\subsection*{A finite counterexample and all assumptions}
Equation reference: walk. Let \(m=2^{17}\) and \(N=m-1=131071\). On \(\Omega=\{-1,1\}^{N}\) with the uniform probability measure let \(\varepsilon_1,\ldots,\varepsilon_N\) be the coordinate signs. Define the absorbed random walk \[
 X_0=1,\qquad X_k=X_{k-1}+\varepsilon_k\mathbf 1_{\{X_{k-1}>0\}}
 \quad(1\le k\le N).\] Take \(\mathcal M=L^\infty(\Omega)\), the probability trace \(\tau(f)=\mathbb Ef\), and the coordinate filtration \(\mathcal M_k\), extended constantly after time \(N\). This is a finite, hence semifinite, von Neumann algebra with a normal faithful trace. The filtration has union \(\mathcal M\) and the ordinary trace-preserving conditional expectations. In particular, no issue concerning an unbounded trace or existence of conditional expectations enters this example.\par
Equation reference: choice. The walk is nonnegative and absorbed at zero. Its increments have zero conditional mean, since \(\mathbf 1_{\{X_{k-1}>0\}}\) is measurable at time \(k-1\). Thus \(X\) is a finite martingale. Set \[
 y_k=-X_k,\qquad x_N=0,\qquad z_N=1,\qquad p=4.\] All variables are real, self-adjoint, and bounded by \(N+1\) in absolute value; in particular every required \(L^p\) membership holds.\par
For any \(\mu>0\), \(y_k/\mu\le0<1\). The projections generated by this scaled martingale are therefore all \(I\). Both sides of [test1] (\#test1) are zero. Moreover, \[(dy_0)^2=1,\qquad
 (dy_k)^2=\mathbf 1_{\{X_{k-1}>0\}}\le1\quad(k\ge1).\] For every allowed projection \(P\), \[\tau\bigl(P(dy_k/\mu)^2P\bigr)
 \le\mu^{-2}\tau(P)
 =\tau\bigl(P(z_N/\mu)^2P\bigr).\] This proves [test2] (\#test2), including \(k=0\), for every positive scale. The large but finite value of \(N\) is admissible: the asserted constant has no dependence on the number of martingale steps or on the algebra dimension.\par
\subsection*{Exact evaluation and contradiction}
Equation reference: reflection. Let \(S_r\) denote an unabsorbed symmetric random walk started at zero. For \(j>0\) of the appropriate parity, reflecting a path after its first visit to zero gives \[\begin{aligned}
 \mathbb P(X_N=j)
 &=2^{-N}\left\{
 \binom{N}{(N+j-1)/2}-\binom{N}{(N+j+1)/2}\right\}\notag\\
 &=\frac{2j}{N+1}\mathbb P(S_{N+1}=j).\end{aligned}\] The first binomial counts all paths from \(1\) to \(j\); the second counts the paths that visit zero, by their bijection with paths from \(-1\) to \(j\). An out-of-range binomial coefficient is zero. The second equality follows directly from the two adjacent binomial coefficients. One may define \(S_{N+1}\) on an auxiliary fair-sign space; only its distribution is used here.\par
Equation reference: lower. Symmetry and [reflection] (\#reflection) give the exact identity \[\mathbb EX_N^4=\frac{\mathbb E|S_{N+1}|^5}{N+1}.\] Since \(\mathbb ES_{N+1}^2=N+1\), Jensen’s inequality for \(u^{5/2}\) yields \[
 \|y_N\|_4=(\mathbb EX_N^4)^{1/4}
 \ge (N+1)^{3/8}=2^{51/8}>80.\] The strict comparison is entirely integral: \(2^{51}>80^8\) is equivalent to \(2^{19}>5^8\), namely \(524288>390625\).\par
In contrast, the right-hand side of [asserted] (\#asserted) is exactly \[\frac{48}{\sqrt{1-(5/4)^{-2}}}
 (0^2+1^2)^{1/2}=80.\] Hence [lower] (\#lower) contradicts the stated conclusion. This is an analytic finite counterexample, not a simulation or an extrapolation. The independently written calculation file checks the constant and integer comparison exactly and verifies [reflection] (\#reflection) against direct finite path counting for small lengths; those finite checks support, but do not replace, the argument above.\par
\subsection*{The precise failed inference and counterevidence}
The missing sign condition fails explicitly in this same construction. For \(-y=X\) at \(\mu=1\), the strict cutoff gives \(R_0(X)=0\) because \(X_0=1\). The \(n=0\) contribution to the left side of [test1] (\#test1) is then \[\sum_{k=1}^{N}\mathbb E(dX_k)^2\ge \mathbb E(dX_1)^2=1,\] whereas its right side is zero. Thus testing for \(-y\) is not an automatic consequence of testing at every positive scale for \(y\).\par
There is no zero-initial-value assumption to exclude the construction. Self-adjointness does not impose nonnegativity. The quantifier \(\mu>0\) does not include a change of sign. Definition 3.3 would provide stronger information, but cannot silently replace Definition 3.2 in the statement being tested. The actual page images confirm these formulas, excluding a text-extraction error as the explanation.\par
The one-sided Theorems 3.4 and 3.5 are not contradicted: their tail projections are trivial for the chosen nonpositive martingale. The assertion of \(L^p\) membership in Theorem 3.6 is not separately refuted; all counterexample variables are bounded.\par
\section*{Impact and a bounded repair}
The supported assessment is Major for the acquired author version. Theorem 3.6 is explicitly presented as a main result and supplies the abstract passage from testing conditions to moment estimates. Its norm conclusion fails, with every stated assumption satisfied. A short possible repair does not reduce the mathematical impact of that false core implication.\par
A precise necessary change to the proof strategy is to require the Definition 3.2 tests for both \[(x_N/\mu,y/\mu,z_N/\mu)
 \quad\text{and}\quad
 (x_N/\mu,-y/\mu,z_N/\mu)
 \qquad\text{for every }\mu>0.\] These assumptions supply the missing input for (3.29). In the finite-dimensional tracial setting the inspected one-sided arguments, followed separately for the two signs and combined as on p. 19, give [asserted] (\#asserted) with the same constant. This verifies a restricted repair on that domain. A new, independent proof of every unbounded semifinite technical passage is not claimed. The sign-invariant strong conditions in Definition 3.3 also remove the demonstrated obstruction.\par
Theorem 4.1 provides concrete counterevidence against treating all applications as false. Its choice \(x_N=z_N=(\sum_kdy_k^2)^{1/2}\) and its verification on p. 20 work for either sign of \(y\). This application survives the demonstrated sign correction. The other downstream applications require their own sign checks and are outside the present verified scope. No experiments or author-code execution are required to establish the counterexample.\par
A bounded search for an erratum or corrigendum and inspection of the arXiv and publisher landing pages identified no correction to this exact statement. This does not establish priority or the absence of a correction. The later publisher abstract differs from the author manuscript, reinforcing the need for a separate publisher-text comparison before any journal-version attribution.\par
\section*{Conclusion}
Dimension | Recorded state\par
Issue judgment | Confirmed false norm bound; Major in arXiv v2.\par
Repair verification | Both-sign input closes the identified step; same bound checked in the finite-dimensional tracial setting.\par
Scientific decision | Pending; initial review is not independent acceptance.\par
Author adoption | Unknown; no author contact or adoption claim.\par
Remaining work | Publisher-text comparison; unexamined applications and general semifinite repair details.\par
\section*{Supporting checks}
The supplied PDF has SHA256\par
`f824e33da73ab98f6045f7beae5a9a890`
`e6e3e457675b4ecc5f4c5e567df1970`.\par
The source-acquisition searches did not identify a separate supplement or computational repository; this absence is not classified as a mathematical defect.\par
No human review is recorded and no author code was executed. The earlier candidate checkpoint and rendering-time false inspection records remain preserved alongside the later completed checks.\par
9 Y. Jiao, A. Osękowski, and L. Wu, Noncommutative Good-\(\lambda\) Inequalities, arXiv:1805.07057v2, 19 November 2018, 52 pp. https://arxiv.org/abs/1805.07057v2. EMS Press, Noncommutative good-\(\lambda\) inequalities, article landing page for DOI 10.4171/JEMS/1644, publication dated 7 May 2025, consulted 27 October 2026 UTC. https://ems.press/journals/jems/articles/14298745.\par
\end{document}
